% !TEX root = ../main.tex

% The chapter letters distinguish the two halves of Talk 4. Inside both
% halves, sections and numbered statements use the same 4.x sequences.
\renewcommand{\thechapter}{\arabic{chapter}a}
\chapter{The ordinary finite-limit completion of manifolds}
\label[chapter]{chap:Ordinary_Finite_Limit_Completion_Manifolds}
\renewcommand{\thechapter}{\arabic{chapter}}

\emph{Seminar note. This chapter and \Cref{chap:From_Smooth_Algebra_To_Derived_Manifolds}
were covered together in Talk~4.}

In \Cref{thm:Manifolds_Inside_Smooth_Algebra} we embedded manifolds in
the opposite category of finitely presented $\Cinfty$-rings. There
every pair of maps to a common target has a pullback, and transverse
maps of manifolds have the pullback we already know from differential
geometry. Could we have arrived at this category by asking for just
these properties?

\begin{question}
  \label[question]{question:Universal_Finite_Limit_Completion}
  What is the universal way to adjoin finite limits to the category of
  manifolds while preserving the point and transverse pullbacks?
\end{question}

Write $\operatorname{Aff}_{\Cinfty}^{\mathrm{fp}}$ for the opposite
category of finitely presented $\Cinfty$-rings and
$i_0\colon\mathrm{Mfd}\to\operatorname{Aff}_{\Cinfty}^{\mathrm{fp}}$
for the functor taking $M$ to $\Cinfty(M)$. This category has finite
limits because finitely presented rings are closed under finite colimits.
Concretely, a coproduct combines the generators and relations in two
presentations, and a coequalizer imposes equality on the images of
each of the finitely many generators of its source.

\begin{theorem}[The ordinary finite-limit completion of manifolds]
  \label[theorem]{thm:Ordinary_Finite_Limit_Completion_Manifolds}
  For every ordinary category $C$ with finite limits, restriction along
  $i_0$ induces an equivalence of categories
  \[
    \Fun^{\mathrm{lex}}(\operatorname{Aff}_{\Cinfty}^{\mathrm{fp}},C)
    \simeq \Fun^{\pitchfork}(\mathrm{Mfd},C).
  \]
  On the left are the functors preserving finite limits. On the right
  are the functors preserving the point and transverse pullbacks.
\end{theorem}

We will construct the extension of a functor on manifolds by solving
finite systems of smooth equations in $C$. We then ask what these
equations tell us about a tangency and about a self-intersection.

\section{Extending functors by generators and relations}
\label[section]{sec:Ordinary_Finite_Limit_Envelope}

A functor on manifolds preserving the point and transverse pullbacks
preserves products, since a product is a transverse pullback over the
point. Its restriction to the Euclidean spaces therefore has the
following structure.

\begin{definition}[$\Cinfty$-ring object]
  \label[definition]{def:Internal_Cinfty_Ring}
  Let $C$ be an ordinary category with finite products. A \emph{$\Cinfty$-ring object in
  $C$} is a finite-product-preserving functor
  \[
    R\colon\mathrm{Euc}\longrightarrow C.
  \]
  We write $\operatorname{Alg}_{\Cinfty}(C)$ for the category of
  $\Cinfty$-ring objects in $C$.
\end{definition}

Let $j\colon\mathrm{Euc}\to\operatorname{Aff}_{\Cinfty}^{\mathrm{fp}}$
be the restriction of $i_0$, taking $\R^n$ to the free ring on $n$ generators.

\begin{theorem}[Ordinary finite-limit envelope]
  \label[theorem]{thm:Ordinary_Cinfty_Finite_Limit_Envelope}
  Let $C$ be an ordinary category with finite limits. Restriction along $j$
  induces an equivalence of categories
  \[
    \Fun^{\mathrm{lex}}
    \bigl(\operatorname{Aff}_{\Cinfty}^{\mathrm{fp}},C\bigr)
    \simeq
    \operatorname{Alg}_{\Cinfty}(C),
  \]
  where the left side consists of finite-limit-preserving functors.
\end{theorem}

\begin{proof}
  Let $R\colon\mathrm{Euc}\to C$ preserve finite products. Given a
  presentation $A=\Cinfty(\R^n)/(f_1,\ldots,f_q)$, assemble the
  equations into $f\colon\R^n\to\R^q$. Define $\widehat R(A)$ by the pullback
  \[
  \begin{tikzcd}
    \widehat R(A) \rar \dar & R(\R^n) \dar["R(f)"] \\
    R(*) \rar["R(0)"'] & R(\R^q).
  \end{tikzcd}
  \]
  Here $0$ denotes the origin and $R(*)$ is terminal. When $A$ is free
  there are no equations, so this construction returns $R(\R^n)$.

  To see why the presentation can be chosen freely, fix $c\in C$.
  The sets $B_c(T)=\Hom_C(c,R(T))$ form an ordinary $\Cinfty$-ring
  as $T$ ranges over Euclidean spaces. A map $c\to\widehat R(A)$
  amounts to an $n$-tuple in $B_c(\R)$ solving the equations $f=0$,
  or equivalently a map of $\Cinfty$-rings $A\to B_c$. Thus
  \[
    \Hom_C(c,\widehat R(A))\cong
    \Hom_{\Cinfty\text{-}\mathrm{Ring}}(A,B_c)
  \]
  naturally in $c$. The right side depends only on $A$ and is
  contravariantly functorial in $A$. By Yoneda this identifies
  the objects obtained from different presentations and defines all
  maps of the extension. It also proves left exactness, since mapping
  a colimit of rings into $B_c$ gives a limit of sets.

  A natural transformation $R\to R'$ acts on these solution objects
  by the universal property of the displayed pullback. Conversely,
  every finitely presented $A$, viewed in the opposite category, is
  itself the displayed zero fibre for $j$. A left-exact extension is
  therefore determined on objects by its restriction to Euclidean
  spaces. For a map of presented rings, choose smooth representatives
  for the images of the generators. Applying $R$ to this tuple induces
  the map on solution objects, uniquely determined by substitution on
  solutions. The same pullback determines natural
  transformations from their restrictions. This proves the equivalence.
\end{proof}

For the ordinary envelope of an algebraic theory see also
\citet[Remark~3.26]{Carchedi_Steffens_Universal_Derived_Manifolds}.

\begin{proof}[Proof of \Cref{thm:Ordinary_Finite_Limit_Completion_Manifolds}]
  Restriction along $i_0$ is well defined by
  \Cref{thm:Manifolds_Inside_Smooth_Algebra} and $\Cinfty(*)=\R$.
  Given $F\in\Fun^{\pitchfork}(\mathrm{Mfd},C)$, write $R$ for its
  restriction to Euclidean spaces and form $\widehat R$ as above.
  We must identify $F(M)$ with $\widehat R(\Cinfty(M))$ naturally in $M$.

  There is a comparison map before we choose any presentation of $M$.
  For $u\colon c\to F(M)$, send a smooth function $h\colon M\to\R$ to
  $F(h)u\in B_c(\R)$. This assignment respects every smooth operation
  since $F$ preserves products. It defines a ring map $\Cinfty(M)\to B_c$,
  hence a map $c\to\widehat R(\Cinfty(M))$. By Yoneda we obtain the
  desired comparison, natural in $M$ and $F$. On Euclidean spaces it
  is an isomorphism by construction.

  For an open $U\subseteq\R^d$, \Cref{prop:Open_Subset_One_Relation}
  presents $U$ as the regular zero locus of $a(x)y-1$ in $\R^{d+1}$.
  Both functors preserve this transverse pullback, so the comparison
  is an isomorphism on $U$. Every manifold is a retract of an open
  tubular neighbourhood in Euclidean space. By naturality, the
  comparison on a manifold is a retract of an isomorphism and hence
  an isomorphism. In the other direction, extending the restriction
  of a left-exact functor recovers that functor by
  \Cref{thm:Ordinary_Cinfty_Finite_Limit_Envelope}.
\end{proof}

Ordinary categories with finite limits also split idempotents. Indeed,
if $e\colon X\to X$ satisfies $e^2=e$, take the equalizer
$i\colon Y\to X$ of $e$ and $\operatorname{id}_X$. Then $e=ir$ for
a unique $r\colon X\to Y$, and $iri=ei=i$ implies $ri=\operatorname{id}_Y$
since $i$ is monic.

\section{Two nontransverse intersections}
\label[section]{sec:What_Ordinary_Envelope_Forgets}

For a smooth map $f\colon M\to\R^q$, the zero locus in the ordinary
envelope thus has ring $\Cinfty(M)/(f_1,\ldots,f_q)$.

\begin{example}[A transverse zero]
  \label[example]{ex:Ordinary_Envelope_Transverse_Zero}
  The quotient $\Cinfty(\R)/(x)\cong\R$ says that the formal zero locus
  of $x$ is the ordinary point, as required by transverse compatibility.
\end{example}

\begin{example}[A singular zero]
  \label[example]{ex:Ordinary_Envelope_Singular_Zero}
  The dual-number quotient $\Cinfty(\R)/(x^2)$ from the preceding
  chapters is the formal zero locus of $x^2$. Thus the ordinary envelope
  already retains its nilpotent thickening.
\end{example}

Now consider the self-intersection of a closed subobject. Let $A$ be a
$\Cinfty$-ring, let $I\subseteq A$ be an ideal, and put
\[
  B=A/I.
\]
The self-intersection of the corresponding formal closed subobject is computed
by
\[
  B\amalg_A B.
\]

\begin{proposition}[Ordinary self-intersections]
  \label[proposition]{prop:Ordinary_Self_Intersections_Collapse}
  There is a natural isomorphism
  \[
    (A/I)\amalg_A(A/I)
    \cong
    A/I.
  \]
\end{proposition}

\begin{proof}
  Apply the quotient pushout formula:
  \[
    (A/I)\amalg_A(A/I)
    \cong
    A/(I+I)
    =A/I.
  \]
\end{proof}

In particular, taking $A=\Cinfty(\R)$ and $I=(x)$ shows that the
ordinary self-intersection of the origin has ring $\R$. The two
nontransverse examples give different answers. In the tangency, the
class of $x$ survives and has square zero. In the self-intersection,
the pullback is again the point. Where did the excess direction from
\Cref{ex:Excess_Self_Intersection} go?
